\section{Methodology}\label{sec:methodology}

\subsection{Inference under the physical measure}

The baseline statistical model assumes that the observed underlying follows a geometric Brownian motion under the physical measure,
\begin{equation}
 dS_t=\mu S_t\,dt+\sigma S_t\,dW_t^{P}.
 \label{eq:p_gbm}
\end{equation}
For equally spaced observations with interval $\Delta t$, log returns satisfy
\begin{equation}
 R_i=\log\!\left(\frac{S_{t_i}}{S_{t_{i-1}}}\right)
 \sim
 \mathcal N\!\left[
 \left(\mu-\frac{1}{2}\sigma^2\right)\Delta t,
 \sigma^2\Delta t
 \right].
 \label{eq:return_likelihood}
\end{equation}
The likelihood is the product of the conditional return densities. The diffusion is intentionally simple: the paper's primary object is the incremental effect of parameter uncertainty and its interaction with the pricing map, not a claim that constant-volatility GBM is a complete model of crude-oil returns.

The baseline prior specification is
\begin{align}
 \mu&\sim\mathcal N(0,1),\\
 \sigma&\sim\operatorname{InvGamma}(2,0.1),\qquad \sigma>0.
 \label{eq:priors}
\end{align}
Sampling is performed in $(\mu,\eta)$ with $\eta=\log\sigma$ so that positivity is automatic. The transformed log posterior is
\begin{equation}
 \log \pi(\mu,\eta\mid\mathcal D)
 =\log L(\mu,e^\eta;\mathcal D)
 +\log p(\mu)+\log p(e^\eta)+\eta+C,
 \label{eq:posterior}
\end{equation}
where the final $\eta$ is the Jacobian term.

A symmetric random-walk Metropolis proposal is used,
\begin{equation}
 \theta'=\theta+\varepsilon,
 \qquad
 \varepsilon\sim\mathcal N(0,\Sigma_q),
\end{equation}
with $\theta=(\mu,\eta)$. The acceptance probability is
\begin{equation}
 \alpha(\theta,\theta')=
 \min\left\{1,
 \frac{\pi(\theta'\mid\mathcal D)}{\pi(\theta\mid\mathcal D)}
 \right\},
\end{equation}
following the Metropolis--Hastings construction \cite{metropolis1953,hastings1970}. The empirical implementation uses multiple chains and Gelman--Rubin $\hat R$ diagnostics for $\mu$ and $\sigma$. For the large synthetic mechanism and calibration experiments, $\mu$ is also integrated analytically under its Gaussian prior, leaving a one-dimensional posterior in $\sigma$ that is evaluated by dense quadrature. This provides an MCMC-independent numerical reference without changing the statistical model.

\subsection{Risk-neutral valuation}

Historical inference and derivative valuation are kept conceptually and computationally separate. Under the Black--Scholes--Merton benchmark,
\begin{equation}
 dS_t=(r-q)S_t\,dt+\sigma S_t\,dW_t^{Q}.
 \label{eq:q_gbm}
\end{equation}
The historical drift $\mu$ in Equation~\eqref{eq:p_gbm} does not enter Equation~\eqref{eq:q_gbm}. Posterior uncertainty about $\mu$ is therefore relevant for physical forecasting but is not propagated into the no-arbitrage price in the baseline model.

For a discretely monitored arithmetic-average option with fixing dates $t_1,\ldots,t_M$, conditional risk-neutral value is
\begin{equation}
 C^Q(\sigma)=e^{-r\tau}
 \mathbb E_{\sigma}^{Q}\!\left[\Pi(A_T,K)\mid\mathcal F_t\right],
 \label{eq:conditional_price}
\end{equation}
where $\Pi(A_T,K)$ is the call or put payoff and $\tau=T-t$. The posterior for $\sigma$ induces a posterior distribution over theoretical prices through the push-forward
\begin{equation}
 p(C\mid\mathcal D)=
 \int
 \delta_{C^Q(\sigma)}(C)
 p(\sigma\mid\mathcal D)\,d\sigma.
 \label{eq:pushforward}
\end{equation}
This distribution quantifies parameter uncertainty in a model price. It is not a predictive distribution of portfolio losses and is not interpreted as VaR or expected shortfall.

\subsection{WTI futures representation}

For the WTI APO application, the state variables are the first-nearby futures fixings in Equation~\eqref{eq:apo_average}. A parsimonious risk-neutral benchmark for a futures price is
\begin{equation}
 dF_u=\sigma F_u\,dW_u^Q.
 \label{eq:futures_q}
\end{equation}
The implementation does not replace the entire APO with one fixed CL contract. For each remaining fixing date it identifies the first-nearby contract applicable to that date and initializes the stochastic fixing from the corresponding point of the observed futures term structure.

For a valuation date inside the averaging month, simulation starts from Equation~\eqref{eq:partial_fixing}: realized fixings enter deterministically and only the remaining first-nearby settlements are stochastic. Two options with the same strike and final settlement date can therefore have different effective volatility exposure when their realized-average components differ.

The single-factor specification is deliberately interpretable. Under a constant-volatility GBM, the diffusion coefficient does not have to change mechanically when moving from $\mathbb P$ to $\mathbb Q$. Using a historical estimate of $\sigma$ in Equation~\eqref{eq:futures_q} is therefore coherent under the maintained model. It is not asserted to be a complete commodity model: mean reversion, stochastic convenience yield, stochastic rates, and jump or stochastic-volatility dynamics are well-established alternatives for commodity futures and Asian options \cite{schwartz1997,ewaldwu2023,shirayatakahashi2011}. A market-implied effective volatility can therefore differ because observed option prices reflect risk premia and because the one-factor constant-volatility benchmark may absorb omitted stochastic volatility, jumps, smile effects, or richer futures-curve dynamics. This distinction follows the broader option-pricing evidence that changing model specification can be a first-order empirical margin and that additional complexity must be judged by incremental performance rather than realism alone \cite{bakshicaochen1997,cumminsesposito2025}.

\subsection{Monte Carlo and conditioning-based valuation}

Arithmetic-average options generally lack the simple closed form available for geometric averages. The baseline empirical prices are therefore evaluated by Monte Carlo with common random numbers across nearby volatility values. In the standard Asian benchmark, antithetic sampling and a geometric-average control variate follow the logic of \cite{kemnavorst1990}. More generally, variance-reduction controls are used only when their expectation is known under the exact state representation being simulated.

For robustness, the same one-factor risk-neutral dynamics are also priced using the geometric-conditioning approximation of \cite{curran1994}. Curran pricing is not treated as the true exchange implementation; it is an independent numerical benchmark under the paper's maintained $Q$ dynamics. Difficult empirical contracts are further repriced using multi-million-path pseudo-random Monte Carlo and independently scrambled Sobol sequences. Agreement across these engines is used to determine whether the small posterior-integration effects are numerically resolved.

\subsection{Competing valuation rules}

Four baseline estimators are compared:
\begin{align}
 \widehat C_{PI}
 &=\mathbb E[C^Q(\sigma)\mid\mathcal D],\\
 \widehat C_{PM}
 &=C^Q(\mathbb E[\sigma\mid\mathcal D]),\\
 \widehat C_{Mode}
 &=C^Q(\widehat\sigma_{Mode}),\\
 \widehat C_{MLE}
 &=C^Q(\widehat\sigma_{MLE}).
 \label{eq:estimators}
\end{align}
The first is the posterior-integrated rule, sometimes called full-Bayes pricing. The second collapses the posterior to its mean before pricing. The third uses the marginal posterior mode of $\sigma$, and the fourth is a non-Bayesian likelihood benchmark. The marginal volatility mode is not described as a joint MAP estimator.

The main within-model diagnostic is
\begin{equation}
 \Delta_{PI,PM}
 =\widehat C_{PI}-\widehat C_{PM}.
 \label{eq:fb_pm_gap}
\end{equation}
A second-order expansion gives
\begin{equation}
 \Delta_{PI,PM}
 \approx
 \frac{1}{2}C^{Q\prime\prime}(\bar\sigma)
 \operatorname{Var}(\sigma\mid\mathcal D),
 \qquad
 \bar\sigma=\mathbb E[\sigma\mid\mathcal D].
 \label{eq:taylor_gap}
\end{equation}
Equation~\eqref{eq:taylor_gap} is the central mechanism: posterior integration matters only to the extent that posterior dispersion interacts with nonlinear exposure to volatility. The second derivative $C^{Q\prime\prime}(\sigma)$ is the option's volatility curvature, commonly termed \emph{volga} or \emph{vomma}; it need not have the same sign in every contract state. More formally, if $C^Q$ is three times differentiable over the relevant posterior support and $|C^{Q\prime\prime\prime}(\sigma)|\leq M_3$, the Taylor remainder satisfies
\begin{equation}
 \left|
 \mathbb E[R_3\mid\mathcal D]
 \right|
 \leq
 \frac{M_3}{6}
 \mathbb E\!\left[
 |\sigma-\bar\sigma|^3\mid\mathcal D
 \right].
 \label{eq:taylor_remainder_bound}
\end{equation}
The second-order approximation is therefore most reliable when the posterior is concentrated and higher-order curvature is controlled.

This mean-price correction is distinct from the dispersion of the posterior distribution of model prices. A first-order delta-method approximation gives
\begin{equation}
 \operatorname{SD}\!\left(C^Q(\sigma)\mid\mathcal D\right)
 \approx
 \left|C^{Q\prime}(\bar\sigma)\right|
 \sqrt{\operatorname{Var}(\sigma\mid\mathcal D)}.
 \label{eq:price_uncertainty_delta}
\end{equation}
Thus the PI--PM correction is locally of order $\operatorname{Var}(\sigma\mid\mathcal D)$, whereas model-price uncertainty is locally of order $\sqrt{\operatorname{Var}(\sigma\mid\mathcal D)}$. Close PI and PM point prices therefore do not imply a narrow posterior distribution of theoretical prices.

As a plug-in sensitivity check, the paper also considers the posterior root-mean-square volatility
\begin{equation}
 \sigma_{RMS}
 =
 \sqrt{\mathbb E[\sigma^2\mid\mathcal D]}
 =
 \sqrt{\bar\sigma^2+\operatorname{Var}(\sigma\mid\mathcal D)},
 \label{eq:sigma_rms}
\end{equation}
with price $C^Q(\sigma_{RMS})$. The volatility displacement can be written
\begin{equation}
 \sigma_{RMS}-\bar\sigma
 =
 \frac{\operatorname{Var}(\sigma\mid\mathcal D)}
 {\sigma_{RMS}+\bar\sigma}
 \approx
 \frac{\operatorname{Var}(\sigma\mid\mathcal D)}
 {2\bar\sigma},
 \label{eq:sigma_rms_shift}
\end{equation}
so the associated plug-in price change is also locally of order $\operatorname{Var}(\sigma\mid\mathcal D)$. This comparator is not a replacement for posterior integration; it tests how much the near-equivalence depends on using $\mathbb E[\sigma\mid\mathcal D]$ rather than a variance-matching scalar plug-in.

\subsection{APO-implied effective risk-neutral volatility}

To separate posterior integration from volatility specification, the paper also asks what volatility parameter the observed APO settlement requires under the same one-factor pricing map. For contract $i$ on date $t$, the contract-level implied value $\sigma^{APO}_{Q,i,t}$ solves
\begin{equation}
 C^Q_{i,t}(\sigma^{APO}_{Q,i,t})=P^{mkt}_{i,t},
 \label{eq:apo_iv}
\end{equation}
when the settlement lies within the monotone model price map. Cross-sectional common-volatility estimates minimize settlement squared error across contracts on a date, while smile specifications allow implied volatility to vary with moneyness and option type.

These quantities are best interpreted as \emph{Curran-equivalent APO-implied effective volatilities under the maintained one-factor model}. They are inferred from the APO cross-section itself and are therefore not an independent external $Q$-volatility source. Their role is diagnostic and predictive: they measure how far the volatility parameter needed by the option market lies from the historical constant-volatility estimate, and whether that information transfers to later valuation dates.

\subsection{Independent vanilla-WTI risk-neutral surface}\label{sec:external_lo_method}

The APO-implied exercise above is deliberately informative about volatility specification but remains within the target option family. To obtain an external risk-neutral state, the paper also uses standard monthly American WTI options on futures (CME product LO). For vanilla contract $j$ observed on trading reference date $s$, let $F_{j,s}$ be the matching CL futures settlement, $K_j$ the strike, $\tau_{j,s}$ time to expiration, and $r_s(\tau_{j,s})$ the date-specific rate proxy. The contract-level implied volatility $\sigma^{LO}_{Q,j,s}$ solves
\begin{equation}
 V^{A}_{LO}\!\left(F_{j,s},K_j,\tau_{j,s},r_s,\sigma^{LO}_{Q,j,s}\right)
 =P^{LO}_{j,s},
 \label{eq:lo_iv}
\end{equation}
where $V^{A}_{LO}$ is evaluated with an American Cox--Ross--Rubinstein futures-option tree and $P^{LO}_{j,s}$ is the official LO settlement. The production inversion uses 160 tree steps, as recorded in the inversion manifest. The futures process is a martingale under $\mathbb Q$ in the maintained one-factor representation, and early exercise is evaluated at every tree node.

The resulting contract-level IVs are fitted separately by underlying futures contract with the same quadratic log-moneyness design used for the APO forward smile, including call/put interactions. For an APO target on date $t$ and strike $K$, each required underlying $u$ therefore supplies a strictly prior implied volatility
\begin{equation}
 \widehat\sigma^{LO}_{Q,u,t-1}(K),
 \qquad s<t,
 \label{eq:lo_prior_guard}
\end{equation}
evaluated at moneyness computed from the target-date CL futures curve. The previous-day specification uses the latest completed LO date before $t$; the expanding specification uses all earlier dates with exponentially decaying recency weights.

The APO pricing engine retains one effective volatility state. If $w_u$ is the fraction of APO fixing dates mapped to underlying futures contract $u$, the strike-specific external state is
\begin{equation}
 \widehat\sigma^{LO,eff}_{Q,t}(K)
 =
 \left[
 \sum_u w_u
 \left\{\widehat\sigma^{LO}_{Q,u,t-1}(K)\right\}^2
 \right]^{1/2}.
 \label{eq:lo_effective_sigma}
\end{equation}
Equation~\eqref{eq:lo_effective_sigma} is an effective scalar transfer rule, not an exact variance identity for the arithmetic-average payoff. In general, if $U_t$ indexes unresolved fixings and $M$ is the total number of fixings,
\begin{equation}
 \operatorname{Var}_{Q}(A_T\mid\mathcal F_t)
 =
 \frac{1}{M^2}
 \sum_{j,k\in U_t}
 \operatorname{Cov}_{Q}\!\left(
 F_{t_j}^{u_j},F_{t_k}^{u_k}\mid\mathcal F_t
 \right),
 \label{eq:average_variance_covariance}
\end{equation}
so fixing times, futures levels, and cross-maturity dependence all matter. The weighted-RMS rule is used only to map maturity-specific vanilla IVs back into the paper's maintained one-factor scalar-volatility pricing engine. As a direct sensitivity check, an alternative scalar is chosen to match the variance of the unresolved arithmetic-average component under the same one-common-factor representation, using the target-date fixing times and futures levels in Equation~\eqref{eq:average_variance_covariance}. A fixing-weighted arithmetic mean is also available as a simpler diagnostic, but is not treated as a variance identity.

The weights are derived from the target contract's actual fixing schedule rather than hard-coded. Predictions outside the prior observed LO moneyness support are clipped to the nearest support boundary and flagged; the main external comparison is repeated on the exact common-support subset with no clipping. No same-day APO implied volatility enters Equations~\eqref{eq:lo_iv}--\eqref{eq:lo_effective_sigma}.

